\documentclass[letterpaper,10pt]{article}

\usepackage[margin=0.45in]{geometry}
\usepackage[hidelinks]{hyperref}
\usepackage{enumitem}
\usepackage{array}
\usepackage{graphicx}

\pagenumbering{gobble}
\setlength{\parindent}{0pt}
\setlength{\parskip}{0pt}

\setlist[itemize]{leftmargin=*, itemsep=2pt, topsep=2pt, parsep=0pt}

\newcommand{\resumesection}[1]{%
    \vspace{6pt}
    \noindent{\large\bfseries #1}\par
    \vspace{2pt}
    \hrule
    \vspace{4pt}
}

\begin{document}

\begin{center}
    {\LARGE \textbf{Jairaj Sharma}}\\[2.5pt]
    \resizebox{\textwidth}{!}{\mbox{US Citizen \textbar\ Email: \href{mailto:sharma.jairajn@gmail.com}{sharma.jairajn@gmail.com} \textbar\ Phone: 603-921-1481 \textbar\ LinkedIn: \href{https://linkedin.com/in/jairajsharma}{linkedin.com/in/jairajsharma} \textbar\ GitHub: \href{https://github.com/jairaj-S}{github.com/jairaj-S}}}
\end{center}

\vspace{4pt}

\resumesection{Education}
\textbf{Virginia Tech} \hfill Blacksburg, VA\\
Bachelor of Science in Computer Science, GPA: 3.5/4.0 \hfill Aug. 2023 -- May 2026\\
\textit{Relevant Coursework:} Operating Systems, Computer Organization, Programming Languages, Introduction to AI, Data Analytics \& Visualization, Data Structures \& Algorithms

\resumesection{Projects}

\textbf{CanvasInsight} \textbar\ \textit{Flask, React, SQLite, OAuth 2.0, Canvas LTI, LLM API, AI Assisted Coding} \hfill Spring 2026
\begin{itemize}
    \item Developed and deployed an instructor-facing Canvas LMS tool to deliver automated analysis for student discussion threads, currently utilized in two Summer 2026 Computer Science Ethics courses.
    \item Leveraged the school’s LLM capabilities to surface student engagement patterns, common misconceptions, and rubric-aligned insights across threads containing up to 200 student replies.
    \item Generated organized participation metrics, cross-thread trend summaries, and actionable teaching recommendations directly from structured and unstructured data to significantly reduce manual instructor analysis.
\end{itemize}

\textbf{Patient Data Retrieval Tool @ VTHacks 13 Winner} \textbar\ \textit{JS, Python, Expo, FastAPI} \hfill Sept. 2025
\begin{itemize}
    \item Built an Expo/React Native EMS app for patient identification that captures facial images and surfaces critical medical records, reducing average EMT response times.
    \item Designed and prototyped a FastAPI backend supporting facial-recognition matching for dispatcher-to-EMT communication.
\end{itemize}

\resumesection{Experience}

\textbf{Undergraduate Research Assistant} \hfill September 2025 -- June 2026\\
\textit{Virginia Tech National Security Institute} \hfill Blacksburg, VA
\begin{itemize}
    \item Engineered 2 programmatic UAS (Unmanned Aerial Systems) scenarios in a Virtual Testing Environment, emulating drone base-incursion events via Python API to support replicable adversarial simulation.
    \item Automated waypoint generation and telemetry capture (position, RF path loss, success metrics) for ground and aerial vehicles, enabling reproducible batch experiments across parameterized configurations.
    \item Developed a parameterized experiment generator using PyYAML and SciPy with Latin Hypercube Sampling; built pandas and Matplotlib dashboards to track and analyze trial outcomes.
\end{itemize}

\vspace{4pt}

\textbf{Student Concept Developer, Hacking for Defense Fall 2025 Cohort} \hfill September 2025 -- December 2025\\
\textit{Virginia Tech} \hfill Blacksburg, VA
\begin{itemize}
    \item Selected as 1 of 5 students to lead concept development for single-pilot multi-drone swarm control, translating mission needs into system requirements and risk constraints.
    \item Led stakeholder discovery with Air Force leadership through 70+ structured interviews on base-incursion scenarios, surfacing key evaluation metrics (pilot workload, link latency) to prioritize solution paths.
    \item Analyzed the command-and-control software stack, comparing ground-station workflows (Betaflight vs.\ MAVLink / QGroundControl) against Federal Aviation Administration and DoD deployability constraints.
\end{itemize}

\vspace{4pt}

\textbf{AI Solutions and Government Bids Intern} \hfill June 2025 -- September 2025\\
\textit{Trinity Bridge Consulting Services} \hfill Nashua, NH
\begin{itemize}
    \item Documented end-to-end workflows for 9 ML prototypes across data collection, model training, and monitoring phases, standardizing documentation for government proposal submissions.
    \item Built a National Health and Nutrition Examination Survey 2021-2023 machine learning pipeline to predict dietary supplement ingredient categories from lab and questionnaire data, using Decision Trees, GridSearchCV cross-validation, class balancing, and confusion-matrix analysis.
\end{itemize}

\vspace{4pt}

\resumesection{Technical Skills and Certifications}
\textbf{Certifications:} AWS Certified AI Practitioner (July 2026), AWS Certified Cloud Practitioner (August 2026) \\
\textbf{Languages:} Python, C, Java, JavaScript, SQL, HTML/CSS \\
\textbf{Frameworks/Tools:} AWS Bedrock, AWS Sagemaker, Visual Studio, FastAPI, Flask, React Native, Expo, Docker, \\
Git/GitHub/Gitlab, Linux, REST APIs, YAML, JSON \\
\textbf{Libraries:} Pandas, NumPy, Matplotlib, SciPy, Scikit-learn \\
\textbf{Specialized:} MAVLink, QGroundControl, OpenAI API

\end{document}
